\documentclass{beamer}

\usepackage{beamerthemeSingapore}
\usepackage[brazil]{babel}
\usepackage[T1]{fontenc}
\usepackage[utf8]{inputenc}
\usepackage{graphicx}
\usepackage{url}
\usepackage{xcolor}
\usepackage{booktabs}
\usepackage{ragged2e}
\usepackage{etoolbox}
\usepackage{minted}
\usepackage{array}

\newcommand{\tcb}[1]{\textcolor{blue}{#1}}
\newcommand{\tcr}[1]{\textcolor{red}{#1}}
\newcommand{\tcg}[1]{\textcolor{green!50!black}{#1}}

\setminted{
  fontsize=\footnotesize,
  breaklines=true,
  breakanywhere=true,
  frame=lines,
  framesep=2mm,
  baselinestretch=1.05,
  tabsize=2
}

\apptocmd{\frame}{}{\justifying}{}

\setbeamertemplate{navigation symbols}{}

% slides de prática e de resposta: fundo e título em cores próprias
\newcommand{\estilopratica}{%
  \setbeamercolor{frametitle}{fg=orange!85!black}%
  \setbeamercolor{background canvas}{bg=orange!8}}
\newcommand{\estiloresposta}{%
  \setbeamercolor{frametitle}{fg=green!45!black}%
  \setbeamercolor{background canvas}{bg=green!6}}

\title{07 - JavaScript: DOM e eventos}
\author{Prof. Alex Kutzke}
\date{Outubro de 2026}

\begin{document}

% ==========================================
\begin{frame}
	\begin{center}
		\large{Setor de Educação Profissional e Tecnológica}\\
		Tecnologia em Análise e Desenvolvimento de Sistemas\\
		\vspace{1cm}
		\small{DS122 - Desenvolvimento de Aplicações Web 1}
	\end{center}
	\titlepage
\end{frame}

% ==========================================
\begin{frame}{Objetivos da aula}
	Ao final desta aula você deve ser capaz de:
	\begin{enumerate}
		\item Descrever o DOM e distinguir o arquivo HTML da página aberta;
		\item Selecionar elementos com \texttt{querySelector} e
		      \texttt{querySelectorAll};
		\item Alterar texto, atributos e classes de um elemento;
		\item Montar o catálogo na tela a partir do array de produtos;
		\item Explicar o risco de \texttt{innerHTML} com dados do usuário;
		\item Registrar ouvintes de evento com \texttt{addEventListener};
		\item Ligar a busca e o filtro da aula anterior aos campos da página;
		\item Validar um formulário e dar o retorno na própria página.
	\end{enumerate}
\end{frame}

% ==========================================
\begin{frame}{Roteiro}
	\tableofcontents
\end{frame}

% ==========================================
\section{Árvore}
% ==========================================

\begin{frame}[fragile]{A página alterada pelo console}
	\texttt{catalogo.html} aberto, \texttt{F12}, aba \textbf{Console}:

	\begin{minted}{javascript}
document.querySelector("h1").textContent = "Loja Alterada";
document.querySelector(".produto").remove();
document.querySelectorAll(".produto").length   // 5
	\end{minted}

	\vspace{0.3cm}
	Título trocado, um cartão a menos, nenhum arquivo salvo.

	\vspace{0.3cm}
	Recarregue: tudo volta ao que estava.
\end{frame}

\begin{frame}{DOM}
	\emph{Document Object Model}, modelo de objetos do documento.

	\vspace{0.3cm}
	O navegador lê o HTML e cria um objeto para cada elemento, com:
	\begin{itemize}
		\item propriedades: texto, atributos, classes;
		\item métodos: acrescentar filho, remover-se, registrar evento.
	\end{itemize}

	\vspace{0.3cm}
	A tela é desenhada a partir desses objetos, e redesenhada quando um deles
	muda.

	\vspace{0.3cm}
	Especificação do WHATWG, o mesmo grupo do HTML. Igual em todos os navegadores
	atuais.
\end{frame}

\begin{frame}[fragile]{Uma árvore de nós}
	\begin{columns}
		\begin{column}{0.5\textwidth}
			\begin{minted}[fontsize=\scriptsize,frame=none]{text}
document
  html
    head
      meta
      meta
      title
      link
    body
      header
        h1
        nav
      main
        section
          h2
          div.grade
            article.produto
            article.produto
			\end{minted}
		\end{column}
		\begin{column}{0.45\textwidth}
			Cada item é um \textbf{nó} (\emph{node}).

			\vspace{0.3cm}
			Filho é o que está dentro. O texto de um elemento também é um nó.

			\vspace{0.3cm}
			É a mesma árvore que o seletor \texttt{.produto h3} percorria na aula
			de CSS.
		\end{column}
	\end{columns}
\end{frame}

\begin{frame}{O arquivo e a página aberta}
	\begin{tabular}{p{0.3\textwidth}>{\raggedright\arraybackslash}p{0.6\textwidth}}
		\toprule
		\texttt{Ctrl+U}       & o arquivo como chegou do servidor. Não muda. \\
		\midrule
		aba \textbf{Elements} & o DOM naquele instante. Mostra cada alteração
		feita pelo JavaScript. \\
		\bottomrule
	\end{tabular}

	\vspace{0.5cm}
	Quando o catálogo for montado pelo JavaScript, os cartões vão estar na tela
	e na aba \textbf{Elements}, e não no código-fonte.
\end{frame}

{\estilopratica
\begin{frame}[fragile]{Prática: o ambiente}
	No seu \emph{fork} de \texttt{ds122-dom-assignment}:

	\begin{enumerate}
		\item copie suas páginas, \texttt{css/}, \texttt{imagens/} e
		      \texttt{js/produtos.js}, ou mova o conteúdo de \texttt{base/} para a
		      raiz;
		\item crie \texttt{js/app.js} com \texttt{formataPreco},
		      \texttt{buscaPorNome} e \texttt{filtraPorFaixa} da tarefa anterior;
		\item ligue \texttt{produtos.js} e \texttt{app.js}, nesta ordem, com
		      \texttt{defer};
		\item no console, confira \texttt{produtos.length}.
	\end{enumerate}

	\vspace{0.3cm}
	Tempo: 5 minutos.
\end{frame}
}

{\estiloresposta
\begin{frame}[fragile]{Resposta: o ambiente}
	\begin{minted}{html}
<link rel="stylesheet" href="css/estilo.css">
<script src="js/produtos.js" defer></script>
<script src="js/app.js" defer></script>
	\end{minted}

	\begin{minted}{javascript}
produtos.length   // 6, com a pasta base/
	\end{minted}

	\vspace{0.2cm}
	Sem erro vermelho no console. Com \texttt{app.js} antes de
	\texttt{produtos.js}, o \texttt{app.js} não enxerga o array.
\end{frame}
}

% ==========================================
\section{Seleção}
% ==========================================

\begin{frame}[fragile]{\texttt{querySelector}}
	O \textbf{primeiro} elemento que casa com um seletor CSS:

	\begin{minted}{javascript}
const titulo = document.querySelector("h1");
const grade = document.querySelector(".grade");
const campoNome = document.querySelector("#nome");
const preco = document.querySelector(".produto .preco");
	\end{minted}

	\vspace{0.2cm}
	Também existe em cada elemento, e aí busca só dentro dele:

	\begin{minted}{javascript}
const cartao = document.querySelector(".produto");
cartao.querySelector("h3");
	\end{minted}
\end{frame}

\begin{frame}[fragile]{\texttt{querySelectorAll}}
	\textbf{Todos} os que casam, numa \texttt{NodeList}:

	\begin{minted}{javascript}
const cartoes = document.querySelectorAll(".produto");

cartoes.length     // 6
cartoes[0]         // o primeiro cartão
cartoes.forEach((cartao) => console.log(cartao));
	\end{minted}

	\vspace{0.2cm}
	Tem \texttt{forEach}. Não tem \texttt{map}, \texttt{filter} nem
	\texttt{reduce}; para usá-los, \texttt{Array.from(cartoes)}.

	\begin{minted}{javascript}
cartoes.classList.add("destaque");
// TypeError: Cannot read properties of undefined (reading 'add')
	\end{minted}
\end{frame}

\begin{frame}[fragile]{\texttt{getElementById}}
	A forma anterior, só para \texttt{id}, e ainda em uso:

	\begin{minted}{javascript}
document.getElementById("nome");     // sem o #
document.querySelector("#nome");     // o mesmo elemento
	\end{minted}

	\vspace{0.3cm}
	Nesta aula, \texttt{querySelector} em todos os casos.
\end{frame}

\begin{frame}[fragile]{Quando nada é encontrado}
	\begin{minted}{javascript}
const botao = document.querySelector("#enviar");  // null
botao.addEventListener("click", enviar);
// TypeError: Cannot read properties of null
//            (reading 'addEventListener')
	\end{minted}

	\vspace{0.2cm}
	O erro aponta a segunda linha. A causa está na primeira:
	\begin{itemize}
		\item \texttt{\#} ou \texttt{.} a menos no seletor;
		\item \texttt{id} escrito diferente do HTML;
		\item elemento que está em outra página;
		\item script sem \texttt{defer}.
	\end{itemize}
\end{frame}

\begin{frame}{Por que o script precisa de \texttt{defer}}
	Sem \texttt{defer}, o script executa quando o navegador encontra a tag, no
	\texttt{<head>}.

	\vspace{0.3cm}
	O \texttt{<body>} ainda não foi lido, e todo \texttt{querySelector} devolve
	\texttt{null}.

	\vspace{0.3cm}
	Com \texttt{defer}, o script espera o documento inteiro entrar na árvore.
\end{frame}

% ==========================================
\section{Alteração}
% ==========================================

\begin{frame}[fragile]{Texto e atributos}
	\begin{minted}{javascript}
const titulo = document.querySelector("main h2");
titulo.textContent = `Produtos (${produtos.length})`;

const imagem = document.querySelector(".produto img");
imagem.src = "imagens/cafe.jpg";
imagem.alt = "Pacote de café em grão de 500 g";

const campoNome = document.querySelector("#nome");
campoNome.value          // o que está digitado, sempre string
campoNome.value = "";    // apaga
	\end{minted}

	\vspace{0.2cm}
	\texttt{textContent} troca todo o conteúdo e não interpreta tags.
\end{frame}

\begin{frame}[fragile]{Classes: \texttt{classList}}
	\begin{minted}{javascript}
const aviso = document.querySelector(".aviso");

aviso.classList.add("oculto");
aviso.classList.remove("oculto");
aviso.classList.toggle("oculto");
aviso.classList.contains("oculto")   // true ou false
	\end{minted}

	\begin{minted}{css}
.oculto {
  display: none;
}
	\end{minted}

	\vspace{0.2cm}
	Mexe só na classe informada.\\
	\texttt{class="produto"} vira \texttt{class="produto destaque"}.
\end{frame}

\begin{frame}[fragile]{\texttt{style}, e por que preferir classe}
	\begin{minted}{javascript}
aviso.style.display = "none";
cartao.style.borderColor = "gold";   // border-color
	\end{minted}

	\vspace{0.2cm}
	Funciona, com desvantagens em relação à classe:
	\begin{itemize}
		\item a aparência fica em dois lugares, CSS e JavaScript;
		\item desfazer exige saber o valor anterior;
		\item media query não alcança o que o script escreveu;
		\item a especificidade é a mais alta da cascata.
	\end{itemize}

	\vspace{0.2cm}
	\texttt{style} para valor calculado na hora, como a largura de uma barra de
	progresso.
\end{frame}

{\estilopratica
\begin{frame}[fragile]{Prática: título e aviso}
	No \texttt{app.js}:

	\begin{enumerate}
		\item o \texttt{<h2>} da lista passa a mostrar \texttt{Produtos (6)}, com
		      o número vindo de \texttt{produtos.length};
		\item a faixa \texttt{.aviso} some, por classe.
	\end{enumerate}

	\vspace{0.3cm}
	Depois compare o \texttt{<h2>} no \texttt{Ctrl+U} e na aba
	\textbf{Elements}.

	\vspace{0.3cm}
	Tempo: 6 minutos.
\end{frame}
}

{\estiloresposta
\begin{frame}[fragile]{Resposta: título e aviso}
	\begin{minted}{javascript}
const titulo = document.querySelector("main h2");
titulo.textContent = `Produtos (${produtos.length})`;

const aviso = document.querySelector(".aviso");
aviso.classList.add("oculto");
	\end{minted}

	\begin{minted}{css}
/* ---------- 8. Classes ligadas pelo JavaScript ---------- */
.oculto {
  display: none;
}
	\end{minted}

	No \texttt{Ctrl+U}, o \texttt{<h2>} continua \texttt{Produtos}.
\end{frame}
}

% ==========================================
\section{Criação}
% ==========================================

\begin{frame}[fragile]{\texttt{createElement} e \texttt{append}}
	\begin{minted}{javascript}
// 1. cria, fora da página
const aviso = document.createElement("p");

// 2. configura
aviso.textContent = "Frete grátis acima de R$ 100,00";
aviso.classList.add("destaque");

// 3. acrescenta como último filho
document.querySelector("main section").append(aviso);
	\end{minted}

	\vspace{0.2cm}
	\texttt{append} aceita vários nós, em ordem. Também: \texttt{prepend},
	\texttt{remove}, e \texttt{appendChild} em código antigo.
\end{frame}

\begin{frame}[fragile]{Uma função que monta um cartão}
	\begin{minted}{javascript}
const criaCartao = (produto) => {
  const cartao = document.createElement("article");
  cartao.classList.add("produto");

  const nome = document.createElement("h3");
  nome.textContent = produto.nome;

  const preco = document.createElement("p");
  preco.classList.add("preco");
  preco.textContent = formataPreco(produto.preco);

  cartao.append(nome, preco);
  return cartao;
};
	\end{minted}

	Devolve o cartão; quem chama decide onde ele entra.
\end{frame}

{\estilopratica
\begin{frame}[fragile]{Prática: um cartão}
	\begin{enumerate}
		\item acrescente o campo \texttt{imagem} aos objetos do
		      \texttt{produtos.js} (em \texttt{base/} ele já existe);
		\item escreva a \texttt{criaCartao} com \texttt{h3}, \texttt{img} com
		      \texttt{alt}, \texttt{p.categoria} e \texttt{p.preco};
		\item acrescente à \texttt{.grade} o cartão de \texttt{produtos[0]}.
	\end{enumerate}

	\vspace{0.3cm}
	Resultado: um sétimo cartão, com o mesmo estilo dos outros.

	\vspace{0.3cm}
	Tempo: 8 minutos.
\end{frame}
}

{\estiloresposta
\begin{frame}[fragile]{Resposta: um cartão}
	\begin{minted}[fontsize=\scriptsize]{javascript}
const criaCartao = (produto) => {
  const cartao = document.createElement("article");
  cartao.classList.add("produto");
  const nome = document.createElement("h3");
  nome.textContent = produto.nome;
  const imagem = document.createElement("img");
  imagem.src = produto.imagem;
  imagem.alt = `Imagem de ${produto.nome}`;
  const categoria = document.createElement("p");
  categoria.classList.add("categoria");
  categoria.textContent = produto.categoria;
  const preco = document.createElement("p");
  preco.classList.add("preco");
  preco.textContent = formataPreco(produto.preco);
  cartao.append(nome, imagem, categoria, preco);
  return cartao;
};

const grade = document.querySelector(".grade");
grade.append(criaCartao(produtos[0]));
	\end{minted}
\end{frame}
}

\begin{frame}[fragile]{De um array para a tela}
	\begin{minted}{javascript}
const renderiza = (lista) => {
  grade.replaceChildren();
  lista.forEach((produto) => grade.append(criaCartao(produto)));
};

renderiza(produtos);
	\end{minted}

	\vspace{0.2cm}
	\texttt{replaceChildren()} esvazia a grade. Sem ele, a segunda chamada
	acrescenta cartões depois dos antigos. Em código antigo:
	\texttt{grade.innerHTML = ""}.

	\vspace{0.2cm}
	A lista vem por parâmetro: a mesma função vai mostrar o resultado da busca.
\end{frame}

\begin{frame}[fragile]{\texttt{innerHTML}}
	\begin{minted}{javascript}
cartao.innerHTML = `
  <h3>${produto.nome}</h3>
  <p class="preco">${formataPreco(produto.preco)}</p>
`;
	\end{minted}

	\vspace{0.2cm}
	Mais curto, e interpreta como HTML tudo o que estiver no \texttt{\$\{\}}.

	\vspace{0.2cm}
	Nome cadastrado com tags vira parte da página de todo visitante. Um
	\texttt{<img onerror="...">} executa JavaScript: \emph{cross-site
	scripting}, XSS.

	\vspace{0.2cm}
	Na disciplina, \texttt{textContent} e \texttt{createElement} para
	dados.\\ \texttt{innerHTML} só para trecho fixo, escrito no código.
\end{frame}

{\estilopratica
\begin{frame}[fragile]{Prática: o catálogo renderizado}
	\begin{enumerate}
		\item apague do \texttt{catalogo.html} todos os
		      \texttt{<article class="produto">};
		\item escreva a \texttt{renderiza} e chame
		      \texttt{renderiza(produtos)};
		\item na \texttt{criaCartao}, desconto de 10\% ou mais recebe a classe
		      \texttt{destaque};
		\item escreva a regra \texttt{.destaque} no CSS.
	\end{enumerate}

	\vspace{0.3cm}
	O destaque condicional é requisito da Entrega 2.

	\vspace{0.3cm}
	Tempo: 10 minutos.
\end{frame}
}

{\estiloresposta
\begin{frame}[fragile]{Resposta: o catálogo renderizado}
	\begin{minted}{html}
<div class="grade"></div>
	\end{minted}

	\begin{minted}{javascript}
// dentro da criaCartao, antes do return
if (produto.desconto >= 0.1) {
  cartao.classList.add("destaque");
}
	\end{minted}

	\begin{minted}{css}
.destaque {
  border: 3px solid #b8860b;
}
	\end{minted}

	Na \texttt{base/}: café, mel e sabonete, com 15\%, 10\% e 20\%.
\end{frame}
}

% ==========================================
\section{Eventos}
% ==========================================

\begin{frame}[fragile]{\texttt{addEventListener}}
	\begin{minted}{javascript}
const botao = document.querySelector("#mostrar-promocoes");

const mostraPromocoes = () => {
  renderiza(produtos.filter((produto) => produto.desconto > 0));
};

botao.addEventListener("click", mostraPromocoes);
	\end{minted}

	\vspace{0.2cm}
	Nome do evento e a função, chamada de ouvinte (\emph{listener}).

	\vspace{0.2cm}
	A linha do \texttt{addEventListener} executa uma vez e só registra. A
	\texttt{mostraPromocoes} executa a cada clique.
\end{frame}

\begin{frame}[fragile]{Passar a função, não o resultado}
	\begin{minted}{javascript}
botao.addEventListener("click", mostraPromocoes);     // certo
botao.addEventListener("click", mostraPromocoes());   // errado
	\end{minted}

	\vspace{0.3cm}
	Com parênteses, a função executa ao carregar a página, e o ouvinte
	registrado é o valor devolvido, \texttt{undefined}.

	\vspace{0.3cm}
	O efeito aparece uma vez sozinho, o clique não faz nada, e o console não
	mostra erro.

	\vspace{0.3cm}
	\texttt{onclick="..."} no HTML também funciona, e mistura comportamento com
	estrutura. Aqui os eventos ficam no script.
\end{frame}

\begin{frame}[fragile]{O objeto do evento}
	\begin{minted}{javascript}
botao.addEventListener("click", (evento) => {
  console.log(evento.type);     // "click"
  console.log(evento.target);   // o elemento clicado
});
	\end{minted}

	\vspace{0.3cm}
	O navegador passa o objeto ao ouvinte. O nome do parâmetro é livre:
	\texttt{event}, \texttt{evt} e \texttt{e} são comuns.
\end{frame}

\begin{frame}[fragile]{\texttt{input} e \texttt{change}}
	\begin{tabular}{l>{\raggedright\arraybackslash}p{0.7\textwidth}}
		\toprule
		\texttt{input}  & a cada alteração do valor: tecla, colagem, seta do
		campo numérico \\
		\texttt{change} & quando a alteração é confirmada: ao sair do campo, ou
		ao escolher no \texttt{<select>} \\
		\bottomrule
	\end{tabular}

	\begin{minted}{javascript}
const campoBusca = document.querySelector("#busca");

campoBusca.addEventListener("input", () => {
  renderiza(buscaPorNome(produtos, campoBusca.value));
});
	\end{minted}

	\texttt{buscaPorNome} é a da aula anterior, sem alteração.
\end{frame}

{\estilopratica
\begin{frame}[fragile]{Prática: busca enquanto digita}
	\begin{enumerate}
		\item acima da grade, um formulário com um campo
		      \texttt{type="search"}, \texttt{id="busca"}, e o seu
		      \texttt{<label>};
		\item ouvinte de \texttt{input} que chama a \texttt{renderiza} com o
		      resultado da \texttt{buscaPorNome}.
	\end{enumerate}

	\vspace{0.3cm}
	Teste com \texttt{de}, e depois apague o termo.

	\vspace{0.3cm}
	Tempo: 8 minutos.
\end{frame}
}

{\estiloresposta
\begin{frame}[fragile]{Resposta: busca enquanto digita}
	\begin{minted}{html}
<form class="filtros" role="search">
  <p>
    <label for="busca">Buscar por nome</label>
    <input type="search" id="busca">
  </p>
</form>
	\end{minted}

	\vspace{0.1cm}
	Com \texttt{de}: chá de hibisco, farinha de milho, sabonete de argila, cesta
	de vime.

	\vspace{0.2cm}
	Campo vazio: \texttt{includes("")} é \texttt{true} para qualquer nome, e todos
	voltam.
\end{frame}
}

\begin{frame}[fragile]{\texttt{value} é sempre string}
	\begin{minted}{javascript}
campoMax.value            // "50", mesmo em type="number"
campoMax.value + 1        // "501"
Number(campoMax.value)    // 50
Number("")                // 0
	\end{minted}

	\vspace{0.3cm}
	Converta antes de fazer conta.

	\vspace{0.3cm}
	Preço máximo vazio convertido direto vira \texttt{0}, e nenhum produto
	passa.
\end{frame}

{\estilopratica
\begin{frame}[fragile]{Prática: faixa de preço}
	\begin{enumerate}
		\item dois campos \texttt{type="number"}, \texttt{preco-min} e
		      \texttt{preco-max}, com \texttt{<label>};
		\item uma função \texttt{aplicaFiltros} que lê os três campos e aplica
		      \texttt{buscaPorNome} e \texttt{filtraPorFaixa};
		\item a mesma função como ouvinte de \texttt{input} nos três campos;
		\item campo vazio não filtra nada.
	\end{enumerate}

	\vspace{0.3cm}
	Teste: \texttt{de} na busca, faixa de 10 a 30.

	\vspace{0.3cm}
	Tempo: 10 minutos.
\end{frame}
}

{\estiloresposta
\begin{frame}[fragile]{Resposta: faixa de preço}
	\begin{minted}[fontsize=\tiny]{javascript}
const campoMin = document.querySelector("#preco-min");
const campoMax = document.querySelector("#preco-max");

const aplicaFiltros = () => {
  const minimo = campoMin.value === "" ? 0 : Number(campoMin.value);
  const maximo = campoMax.value === "" ? Infinity : Number(campoMax.value);

  const porNome = buscaPorNome(produtos, campoBusca.value);
  renderiza(filtraPorFaixa(porNome, minimo, maximo));
};

campoBusca.addEventListener("input", aplicaFiltros);
campoMin.addEventListener("input", aplicaFiltros);
campoMax.addEventListener("input", aplicaFiltros);
	\end{minted}

	\texttt{Infinity} é maior que qualquer número. Resultado do teste: chá de
	hibisco e farinha de milho.
\end{frame}
}

% ==========================================
\section{Formulários}
% ==========================================

\begin{frame}[fragile]{\texttt{submit} e \texttt{preventDefault}}
	Envio do formulário: evento \texttt{submit} no \texttt{<form>}, depois a
	requisição para o \texttt{action}, que carrega outra página.

	\begin{minted}{javascript}
const formulario = document.querySelector("form");

formulario.addEventListener("submit", (evento) => {
  evento.preventDefault();   // cancela o envio
});
	\end{minted}

	\vspace{0.2cm}
	Ouvinte no \texttt{<form>} com \texttt{submit}, e não no botão com
	\texttt{click}: o \texttt{submit} só acontece quando o formulário vai mesmo
	ser enviado.

	\vspace{0.2cm}
	Sem PHP ainda, o contato cancela sempre e responde na própria página.
\end{frame}

\begin{frame}[fragile]{Validação nativa e em JavaScript}
	A nativa (\texttt{required}, \texttt{minlength}, \texttt{type}) executa
	\textbf{antes} do \texttt{submit}. Campo obrigatório vazio: o ouvinte nem é
	chamado.

	\vspace{0.2cm}
	O JavaScript cobre o que os atributos não expressam:

	\begin{minted}{javascript}
"   ".length            // 3: passa por minlength="3"
"   ".trim().length     // 0

if (campoNome.value.trim().length < 3) {
  // recusa
}
	\end{minted}

	Nenhuma das duas protege o servidor. Essa validação vem com o PHP.
\end{frame}

\begin{frame}[fragile]{Mensagem de erro no lugar do campo}
	\begin{minted}{html}
<input type="text" id="nome" name="nome" minlength="3" required>
<span class="erro" id="erro-nome"></span>
	\end{minted}

	\begin{minted}[fontsize=\scriptsize]{javascript}
formulario.addEventListener("submit", (evento) => {
  evento.preventDefault();
  erroNome.textContent = "";     // apaga a anterior

  if (campoNome.value.trim().length < 3) {
    erroNome.textContent = "Informe o nome, com pelo menos 3 letras.";
  }
});
	\end{minted}

	Perto do campo, e não num \texttt{alert}, que bloqueia a página e some ao
	fechar.
\end{frame}

{\estilopratica
\begin{frame}[fragile]{Prática: validação do contato}
	Em \texttt{contato.html}, com \texttt{js/contato.js} ligado:

	\begin{enumerate}
		\item \texttt{preventDefault} no \texttt{submit};
		\item nome com menos de 3 caracteres depois do \texttt{trim}: mensagem
		      abaixo do nome;
		\item mensagem com menos de 10: mensagem abaixo dela;
		\item tudo certo: \texttt{Mensagem enviada com sucesso!} na página, e o
		      formulário limpo.
	\end{enumerate}

	\vspace{0.3cm}
	Teste: três espaços no nome e \texttt{oi} na mensagem.

	\vspace{0.3cm}
	Tempo: 12 minutos.
\end{frame}
}

{\estiloresposta
\begin{frame}[fragile]{Resposta: validação do contato}
	\begin{minted}[fontsize=\tiny]{javascript}
formulario.addEventListener("submit", (evento) => {
  evento.preventDefault();
  erroNome.textContent = "";
  erroMensagem.textContent = "";
  sucesso.classList.add("oculto");
  let valido = true;

  if (campoNome.value.trim().length < 3) {
    erroNome.textContent = "Informe o nome, com pelo menos 3 letras.";
    valido = false;
  }
  if (campoMensagem.value.trim().length < 10) {
    erroMensagem.textContent = "A mensagem precisa ter pelo menos 10 caracteres.";
    valido = false;
  }
  if (valido) {
    sucesso.classList.remove("oculto");
    formulario.reset();
  }
});
	\end{minted}
\end{frame}
}

\begin{frame}{Quando dá errado}
	\small
	\begin{tabular}{>{\raggedright\arraybackslash}p{0.45\textwidth}>{\raggedright\arraybackslash}p{0.45\textwidth}}
		\toprule
		Sintoma & Causa provável \\
		\midrule
		\texttt{Cannot read properties of null} & seletor sem resultado, ou
		script sem \texttt{defer} \\ \midrule
		\texttt{(reading 'add')} numa lista & \texttt{classList} direto na
		\texttt{NodeList} \\ \midrule
		clique não faz nada, efeito ao abrir & função chamada no
		\texttt{addEventListener} \\ \midrule
		soma dá \texttt{"501"} & \texttt{value} sem \texttt{Number} \\ \midrule
		cartões repetidos a cada filtro & contêiner não esvaziado \\ \midrule
		página recarrega no envio & \texttt{preventDefault} ausente \\
		\bottomrule
	\end{tabular}

	\vspace{0.3cm}
	Os quatro últimos não aparecem no console. \texttt{console.log} dentro do
	ouvinte confirma se ele foi chamado, e com que valores.
\end{frame}

\begin{frame}{Tarefa}
	Os exercícios valem nota e compõem o item \emph{Exercícios em sala}.

	\vspace{0.3cm}
	Enunciado no \texttt{README.md} de \texttt{ds122-dom-assignment}; prazo na
	UFPR Virtual.

	\vspace{0.3cm}
	Entrega pelo \emph{fork} no GitLab, individual ou em dupla, sem uso de IA
	generativa para produzir o código.

	\vspace{0.5cm}
	Busca, faixa de preço, destaque e validação são requisitos da Entrega 2. O
	que falta, carregar os produtos de um JSON com \texttt{fetch}, é o assunto do
	próximo encontro.
\end{frame}

\end{document}
